\section{Conclusion}
\label{sec:conclusion}

In this paper, we have proposed the CLIP-Enhanced Semantic Manifold Diffusion (CESM-Diff) framework, a novel generative approach designed to address the challenges of Zero-Shot Fault Diagnosis (ZSFD) in industrial systems. By leveraging the continuous and rich semantic space of the pre-trained CLIP text encoder, our model successfully bridges the generative gap between natural language descriptions and high-dimensional physical sensor signals. Our dual-path diffusion architecture synchronizes the evolution of vibratory features and semantic embeddings, ensuring that the synthesized fault representations are both physically plausible and semantically grounded.

The CESM-Diff framework introduces three primary contributions and structural innovations: (1) The first integration of a dual-path diffusion process with the CLIP semantic manifold for industrial sensing, enabling the capture of complex cross-modal relationships; (2) The Semantic Manifold Interpolation (SMI) mechanism, which models the continuous spectrum of mechanical degradation for accurate zero-shot transfer across evolving failure modes; and (3) The use of Attribute Consistency Loss (ACL) as a structural anchor to prevent semantic drift during the iterative denoising process. Furthermore, our approach demonstrates high practical flexibility, as the trade-off between diagnostic fidelity and inference speed can be effectively tuned for diverse hardware settings through timestep adjustment.

Comprehensive experimental results on the TEP, Hydraulic System, and CWRU Bearing datasets demonstrate that CESM-Diff significantly outperforms current state-of-the-art self-supervised and generative baselines. On the CWRU dataset, our model achieved a Top-1 Accuracy of 92.4\% for unseen fault classes, nearly matching the performance of fully supervised methods. Detailed ablation studies confirmed the vital contributions of the CLIP embedding space and the SMI module, showing that the integration of cross-modal semantic knowledge is essential for effective zero-shot transfer in complex industrial environments. The t-SNE visualizations further provided qualitative evidence that our model learns a well-structured and semantically consistent latent manifold, yielding compact and separable clusters even for previously unseen fault types.

Beyond diagnostic accuracy, our analysis highlighted the practical feasibility of the CESM-Diff framework for real-world scaling and edge deployment. By providing a scalable and robust foundation for autonomous health monitoring, this work enables industrial systems to effectively generalize to novel incidents without historical training data. Future research will explore the extension of semantic manifolds to handle multi-fault scenarios and the integration of prognostic capabilities for remaining useful life (RUL) estimation within the same generative framework. We also plan to investigate more computationally efficient denoising kernels and transformer-based diffusion backbones to further reduce inference latency on low-power edge devices. Ultimately, the CESM-Diff framework provides a reliable and mathematically grounded pathway toward the realization of intelligent and autonomous maintenance in the era of Industry 4.0, where the ability to adapt to unseen failure modes is paramount for reliability and safety. By bridging the gap between natural language and physical signals, we have established a framework that significantly reduces the manual overhead of domain-specific feature engineering in zero-shot contexts.
